\section{Separation from the untwisted product baseline}

The untwisted baseline uses the same \(G_{15}\) and \(G_{60}\), but replaces signed half-flip carriers by identity carriers. Thus the internal fiber edges are unchanged, and each slot edge \(\{t,u\}\) contributes carrier edges of the form \((t,x)\sim(u,x)\). This is exactly the Cartesian product \(G_{15}\square G_{60}\), and it is the untwisted product baseline against which the canonical construction is compared.

\begin{center}
\begin{tabular}{lrr}
\toprule
Invariant & Canonical \(X_\sigma\) & Untwisted baseline \\
\midrule
Diameter & 8 & 9 \\
Radius & 6 & 9 \\
\bottomrule
\end{tabular}
\end{center}

Both diameter and radius separate the canonical graph from the untwisted product baseline.

\begin{center}
\begin{tabular}{lr}
\toprule
Baseline eccentricity & Vertex count \\
\midrule
9 & 900 \\
\bottomrule
\end{tabular}
\end{center}

Thus the canonical object is not merely the product \(G_{15}\square G_{60}\); the ring lives in the signed carrier coupling.

